\subsection{Actual affine dimensions and domination of line components}
\label{subsec:actual-affine-line-components}

\begin{theorem}[Whole prime heights under actual integral going-down]
\label{algebra:integral-normal-height}
\lean{EnvelopeProperty.height_eq_under_of_isIntegral_of_hasGoingDown,
EnvelopeProperty.height_eq_under_of_normal_integral}
\leanok
Let $R\to S$ be an integral extension of Noetherian commutative rings
with genuine going-down. Every prime $P$ of $S$ has exactly the height
of its contraction to $R$. In particular this applies when $S$ is a
domain, the coefficient action is faithful, and $R$ is integrally closed.
\end{theorem}
\begin{proof}
Integral incomparability puts the actual prime $P$ among the minimal
primes over the extension of its contraction. The true quotient prime
correspondence therefore identifies its image in that quotient as a
minimal prime of the quotient ring, of height zero. The genuine
going-down height-addition theorem expresses the entire original height
as the contraction height plus this actual quotient height. Removing
the zero summand gives the equality. A faithful domain algebra over a
normal base satisfies actual going-down; the normal specialization uses
that theorem, rather than a height-equality certificate.
\end{proof}

\begin{theorem}[Actual maximal heights and positive-dimensional prime quotients]
\label{algebra:affine-domain-maximal-height}
\lean{EnvelopeProperty.maximalHeight_mvPolynomialFin,
EnvelopeProperty.exists_constant_maximalHeight_of_finiteType_domain,
EnvelopeProperty.maximalHeight_eq_ringKrullDim_of_finiteType_domain,
EnvelopeProperty.positive_dim_quotient_of_height_lt_of_finiteType_domain}
\leanok
\uses{algebra:integral-normal-height}
For any field $F$, every maximal prime of $F[X_1,\ldots,X_n]$ has
height $n$. Every actual finite-type domain $D$ over $F$ has a common
height for all maximal ideals, equal to its entire Krull dimension.
Consequently a prime $P$ of smaller height than $\dim D$ has
$\dim(D/P)>0$. No algebraic-closure or characteristic hypothesis is used.
\end{theorem}
\begin{proof}
Induct on the actual coordinate count. The zero-variable polynomial
algebra is the field. The genuine polynomial-ring equivalence in the
successor step identifies the maximal prime with one in a one-variable
polynomial ring. Jacobson contraction makes its coefficient prime
maximal, and the actual polynomial height theorem increases that
coefficient height by one.

True Noether normalization constructs an injective finite homomorphism
from an actual polynomial algebra to $D$. That base is normal and the
action is faithful; the preceding height theorem identifies every
maximal height in $D$ with that of its actual maximal contraction.
The genuine characterization of Krull dimension by maximal heights
then makes this common height equal to $\dim D$.

If $P$ has smaller height, it cannot be maximal. If its actual domain
quotient had dimension at most zero, that quotient would be a field.
The true quotient-field criterion would make $P$ maximal, contradicting
the proved strict height inequality. This derives positive dimension
from the actual ring, rather than assuming a component-dimension witness.
\end{proof}

\begin{theorem}[Actual dimension under normal integral extension and principal localization]
\label{algebra:affine-domain-dimension-comparison}
\lean{EnvelopeProperty.ringKrullDim_eq_of_normal_integral_finiteType_domain,
EnvelopeProperty.ringKrullDim_localizationAway_of_finiteType_domain}
\leanok
\uses{algebra:integral-normal-height,algebra:affine-domain-maximal-height}
A faithful integral domain extension between finite-type $F$-algebras
with normal base has the same actual Krull dimension as its base.
For any finite-type domain $R$ over a field and $r\ne0$, the genuine
principal localization $R[1/r]$ also has dimension $\dim R$.
\end{theorem}
\begin{proof}
Choose a genuine maximal ideal in the nontrivial extension domain.
Integral contraction is maximal. The normal-base height equality and
the preceding actual maximal-height identities compute both entire
dimensions from these two actual ideals.

The nonzero principal localization is genuinely a domain and of finite
type over $F$. The original affine ring is Jacobson. Choose an actual
maximal localization ideal; the true Jacobson localization correspondence
makes its contraction maximal. Actual localization preserves this
prime's height. The two maximal-height identities therefore identify
the whole original and localized dimensions. No assertion that all
localizations preserve dimension is used.
\end{proof}

\begin{theorem}[Positive dimension forces actual injectivity over a curve base]
\label{algebra:positive-dimension-integral-injectivity}
\lean{EnvelopeProperty.injective_algebraMap_of_isIntegral_of_positive_dim}
\leanok
Let $C$ be a domain of dimension at most one and $D$ an integral domain
$C$-algebra with actual positive Krull dimension. Then the actual map
$C\to D$ is injective.
\end{theorem}
\begin{proof}
The kernel is prime because $D$ is a domain. A nonzero such kernel
would be maximal in the dimension-at-most-one domain $C$. Integral
maximal-ideal ascent would then make the zero ideal of $D$ maximal,
so $D$ would be a field and have dimension zero. This contradicts its
actual positive dimension; hence the whole original map has zero kernel.
\end{proof}

\begin{theorem}[Original ambient dimension and positive minimal line components]
\label{algebra:line-component-positive-dimension}
\lean{EnvelopeProperty.AffineLine.ringKrullDim_finite_over_principalOpen,
EnvelopeProperty.AffineLine.positive_dim_quotient_of_minimalPrime_lineIdeal}
\leanok
\uses{algebra:affine-domain-dimension-comparison,
algebra:affine-domain-maximal-height,algebra:line-equations-height-bound}
For finitely many coordinates $\sigma$, a genuine finite faithful domain
algebra $B$ over $F[X_i,1/g]$, with $g\ne0$, has dimension
$\#\sigma$. For a chosen parameter index $j$, every actual minimal
prime $P$ of the image of the original transverse line ideal has
$\dim(B/P)>0$ whenever the original ambient domain has dimension at
least $\#\sigma$.
\end{theorem}
\begin{proof}
The genuine nonzero polynomial principal localization is a normal
finite-type domain. Finiteness and the true scalar tower make $B$ a
finite-type field algebra, and faithfulness gives the actual normal
integral dimension comparison. The principal-localization dimension
theorem and genuine polynomial dimension theorem compute its dimension
as the coordinate count.

The original transverse equations have at most $\#\sigma-1$ actual
generators, so the genuine Krull height theorem bounds the height of
every actual minimal prime of their entire image by that number. The
chosen index makes $\sigma$ nonempty, hence this bound is strictly
smaller than the assumed original ambient dimension. The preceding
prime-quotient producer yields actual positive dimension. Minimal
component dimension is a conclusion, not an input certificate.
\end{proof}

\begin{theorem}[Original equations describe the real open-line inverse image]
\label{algebra:localized-line-original-preimage}
\lean{EnvelopeProperty.AffineLine.ker_principalLineHom,
EnvelopeProperty.AffineLine.principalLine_preimage_range,
EnvelopeProperty.AffineLine.principalLineSchemeMap_isPullback}
\leanok
\uses{algebra:line-quotient-whole-polynomial,scheme:principal-open-line-map}
For $b_j\ne0$, the kernel of the actual localized line map
$F[X_i,1/g]\to F[T,1/g(a+bT)]$ is the localization of the original
transverse ideal $J_b$. Under any whole ambient algebra map, its actual
scheme-support inverse image is precisely the zero locus of the entire
original ideal image. The principal-open line square is a genuine
scheme pullback, including empty-open cases.
\end{theorem}
\begin{proof}
Actual localization of a ring homomorphism commutes with its entire
kernel. Apply this to the genuine original substitution and its proved
exact kernel $J_b$. The localized substitution is genuinely surjective,
so its actual prime-spectrum image is the zero locus of this entire
localized kernel. Pulling back ideal containment through the whole
ambient homomorphism and using actual ideal-image composition gives
exactly the zero locus of the original equation image, without radical
replacement in that image ideal.

The whole original localized line square already commutes. Its source
open immersion has precisely the inverse-image open range, by the
genuine localization range formula and actual substitution on $g$.
The true open-immersion pullback theorem therefore makes this square
an actual scheme pullback; no cartesian-square certificate is assumed.
\end{proof}

\begin{definition}[Actual whole open-line map to an original minimal component]
\label{algebra:line-component-whole-map}
\lean{EnvelopeProperty.AffineLine.componentLine_ker_le,
EnvelopeProperty.AffineLine.componentLineHom,
EnvelopeProperty.AffineLine.componentLineHom_comp,
EnvelopeProperty.AffineLine.componentLineHom_principalLineHom}
\leanok
\uses{algebra:localized-line-original-preimage}
Let $\varphi:F[X_i,1/g]\to B$ be an actual $F$-algebra homomorphism
and $P$ an actual minimal prime over the image of $J_b$.
The entire composite to $B/P$ genuinely factors through the surjective
localized line map, constructing an actual homomorphism
\[
 F[T,1/g(a+bT)]\longrightarrow B/P.
\]
Its composite with the original localized substitution equals the
original map to $B/P$ on every localized ambient element.
\end{definition}
\begin{proof}
Minimal-prime membership gives actual containment of the entire
transverse image in $P$. The preceding whole kernel computation and
the actual ideal map/comap adjunction put the kernel of the surjective
localized line map in the kernel of the original quotient composite.
Descend this entire algebra homomorphism by the true surjective-map
universal property. Its whole composite identity and application to
every original localized element give the claimed full coefficient
diagram. Coordinate-value equality is not substituted for this map.
\end{proof}

\begin{theorem}[Genuine finite, injective and dimension-one component maps]
\label{algebra:line-component-domination}
\lean{EnvelopeProperty.AffineLine.componentLineHom_finite,
EnvelopeProperty.AffineLine.componentLineHom_injective,
EnvelopeProperty.AffineLine.componentLineHom_comap_surjective,
EnvelopeProperty.AffineLine.component_quotient_ringKrullDim_eq_one}
\leanok
\uses{algebra:line-component-whole-map,algebra:line-component-positive-dimension,
algebra:positive-dimension-integral-injectivity,
algebra:affine-domain-dimension-comparison,algebra:principal-open-line-ring}
For a finite ambient homomorphism the actual component-line map is
finite. If the ambient map is also injective, $B$ is a domain,
$g\ne0$, and $g(a+bT)\ne0$, then this whole component-line map is
injective, its prime-spectrum map is surjective, and $\dim(B/P)=1$.
No algebraic closure, perfectness or characteristic restriction is used.
\end{theorem}
\begin{proof}
The original quotient $B/P$ is finite over the ambient principal-open
ring. The proved whole coefficient diagram gives the actual scalar
tower through the line ring, so restriction of scalars yields genuine
finiteness over that line ring and hence integrality.

Actual ambient injectivity gives a faithful action. Its computed
dimension and the transverse-equation theorem make $B/P$ genuinely
positive-dimensional. The original prime $P$ makes this quotient a
domain. The nonzero line restriction makes the actual line ring a
Dedekind domain and in particular a domain of dimension at most one.
The positive-dimensional integral-domain producer therefore proves
injectivity of the whole constructed line map. Genuine integral
lying-over now yields surjectivity on the entire prime spectrum.

The actual line base is normal and of finite type, and the actual
component quotient is of finite type. The proved faithful integral
dimension comparison makes its dimension equal to that of this true
nonempty principal-open polynomial line, namely one. Nilpotents in the
original equation quotient are not silently removed; the component
quotient uses the actual original prime $P$.
\end{proof}

\begin{theorem}[Actual finite surjective morphisms of the component schemes]
\label{algebra:line-component-scheme-domination}
\lean{EnvelopeProperty.AffineLine.componentLineSchemeMap,
EnvelopeProperty.AffineLine.componentLineSchemeMap_isFinite,
EnvelopeProperty.AffineLine.componentLineSchemeMap_surjective,
EnvelopeProperty.AffineLine.componentLineSchemeMap_square}
\leanok
\uses{algebra:line-component-domination,algebra:line-component-whole-map}
The genuine scheme morphism
$\operatorname{Spec}(B/P)\to\operatorname{Spec}F[T,1/g(a+bT)]$
is finite and, under the preceding actual domain/injectivity hypotheses,
surjective. Its whole original component inclusion and ambient finite
map commute with the true open-line map.
\end{theorem}
\begin{proof}
Apply the actual Spec functor to the constructed entire line algebra
homomorphism. The genuine affine finiteness criterion transfers the
proved ring finiteness. Its underlying whole prime map is precisely
the already proved surjective comap, giving actual scheme surjectivity.
Contravariant functoriality transfers the entire coefficient composite
identity to the displayed actual scheme square.
\end{proof}

These producers supply the component dimension and domination step in
FHX Proposition~5.22 once the actual finite injective ambient map is
provided. Domain preservation of a scalar-extended ambient algebra,
normalization finiteness and smoothness or finite-flatness of its
normalization are separate genuine algebraic producers. This chapter
assumes neither a component-dominance witness nor a normalization or
curve-potential certificate.
